\section{Exact near-isotropic coefficient}\label{sec:exact-local}
% R01 bounded mathematical investigation, proof_r01, 2026-09-22.
% New proof: do not alter R00 freeze. All implicit constants depend only on
% the fixed word (and its chosen finite period), uniformly over maximizing
% invariant probabilities and their persistent active supports.

\begin{theorem}[Exact cubic coefficient for a fixed BB word]
\label{thm:r01-exact-cubic}
Fix a nonconstant periodic word
$\bm\sigma=(\sigma_0,\ldots,\sigma_{L-1})\in\{0,1\}^L$,
where zero denotes BB1 and one denotes BB2.  Let $(a_j)$ be the unique
$L$-periodic solution of
\[
 a_{j+1}-a_j+2a_{j-1}=2\sigma_j,
\]
and put
\[
 r_j=1-a_{j-1}+\sigma_j,
 \qquad V_{\bm\sigma}=\frac1L\sum_{j=0}^{L-1}(r_j-1)^2.
\]
For $C=(b-a)/(b+a)\downarrow0$, with the word fixed,
\[
 C-\Gamma_{\bm\sigma}(a,b)
 =\frac{V_{\bm\sigma}}2C^3+O_{\bm\sigma}(C^4).
\]
The same expansion holds with $\Gamma_{\bm\sigma}$ replaced by the
finite-dimensional worst-trajectory factor $q^{\rm fd}_{\bm\sigma}$,
the Hilbert worst-trajectory factor $q^{\rm H}_{\bm\sigma}$, or the
finite-dimensional uniform-envelope threshold
$\Theta^{\rm fd}_{\bm\sigma}$, with arbitrary positive warm-up throughout.
In particular,
\[
 \lim_{C\downarrow0}\frac{C-\Gamma_{01}}{C^3}
 =\lim_{C\downarrow0}\frac{C-q^{\rm fd}_{01}}{C^3}
 =\frac1{32}.
\]
\end{theorem}

\begin{proof}
Positive spectral rescaling does not change any factor, so only $C$ matters.
All estimates below are uniform over the maximizing invariant probability
chosen at a given $C$; constants may depend on the fixed word.  The endpoint
response factor is
\[
 \rho_{\bm\sigma}(\kappa)
 =\frac{\kappa-1}
 {\bigl[\prod_j(1+\kappa^{r_j})\bigr]^{1/L}},
 \qquad \kappa=\frac{1+C}{1-C}.
\]
Since $L^{-1}\sum_j r_j=1$, its Taylor expansion gives
\begin{equation}\label{eq:r01-endpoint-loggap}
 \log\frac{C}{\rho_{\bm\sigma}}
 =\frac{V_{\bm\sigma}}2C^2+O_{\bm\sigma}(C^4).
\end{equation}
Indeed, for $z=\log\kappa$, the difference between the period average
of $\log(1+e^{r_jz})$ and $\log(1+e^z)$ is
$V_{\bm\sigma}z^2/8+O_{\bm\sigma}(z^4)$, and
$z=2C+O(C^3)$.

Choose an ergodic maximizing invariant probability $\Pi$, whose existence
and finite logarithmic average follow from
$\Gamma_{\bm\sigma}\ge\rho_{\bm\sigma}>0$ and the attainment theorem.
Let $[A,B]$ be the convex hull of its persistent support, and put
\[
 m=\frac{A+B}{2},\quad d=\frac{B-A}{2},\quad
 r=\frac dm\le C,\quad
 \beta=\int f\,d\Pi,\quad
 G=\log C-\beta,\quad \Delta=\log r-\beta.
\]
The endpoint-band bound and \eqref{eq:r01-endpoint-loggap} imply
\begin{equation}\label{eq:r01-basic-smallness}
 0\le\Delta\le G\le\frac{V_{\bm\sigma}}2C^2+O_{\bm\sigma}(C^4),
 \qquad
 0\le\log(C/r)\le G.
\end{equation}
Consequently $r/C=1+O_{\bm\sigma}(C^2)$, uniformly over $\Pi$.

Write $x=(\lambda-m)/d\in[-1,1]$ and, for a probability $w$, define
\[
 \mu(w)=\int x\,dw,\qquad
 J(w)=1-\int x^2\,dw.
\]
At $X=(w,z,j)$ put $s(X)=(U_j(z)-m)/d$.  The normalized BB selectors are
\[
 s_0(w)=\mu(w),\qquad
 s_1(w)=\frac{\mu(w)+r(1-J(w))}{1+r\mu(w)}.
\]
The endpoint logarithmic defect satisfies the already established bound
\[
 \Delta\ge\mathbb E_\Pi\mathcal D_r(s)
 \ge\frac{(1-C)^2}{2}\mathbb E_\Pi s^2.
\]
Thus $\mathbb E_\Pi s^2=O_{\bm\sigma}(C^2)$.  Moreover,
\[
 |s_1(w)-\mu(w)|
 =\frac{r\bigl(\int x^2\,dw-\mu(w)^2\bigr)}{1+r\mu(w)}
 \le\frac{r}{1-r}.
\]
For the selector actually used at $X$, this gives
$|\mu(z)|\le|s(X)|+r/(1-r)$.  The two coordinate marginals of $\Pi$
coincide, so
\begin{equation}\label{eq:r01-second-moments}
 \mathbb E_\Pi\mu(w)^2
 \le2\mathbb E_\Pi s^2+\frac{2r^2}{(1-r)^2}
 =O_{\bm\sigma}(C^2).
\end{equation}

We next strengthen the interior-mass estimate.  For
$Z=\int(x-s)^2\,dw=1-J-2s\mu+s^2>0$,
\[
 f-\log r=\tfrac12\log Z-\log(1+rs)
 \le-\tfrac12J-s\mu+\tfrac12s^2+\frac{r}{1-r}|s|.
\]
Averaging and applying Cauchy--Schwarz gives
\begin{align*}
 \mathbb E_\Pi J
 &\le2\Delta+2\sqrt{\mathbb E_\Pi s^2\,
                          \mathbb E_\Pi\mu^2}
      +\mathbb E_\Pi s^2
      +\frac{2r}{1-r}\sqrt{\mathbb E_\Pi s^2}\\
 &=O_{\bm\sigma}(C^2).
\end{align*}
This step uses the small second moment of the current mean instead of the
pointwise bound $|\mu|\le1$.

For completeness, the following uniform estimate controls a normalized
update without assuming its normalizer is bounded away from zero:
\begin{equation}\label{eq:r01-mean-remainder}
 \left|\mu(F_uw)-\mu(w)+2s\right|
 \le8\bigl(J(w)+\mu(w)^2\bigr)+12s^2,
 \quad u=m+ds.
\end{equation}
To prove it, let $m_3=\int x^3\,dw$.  The numerator of the left-hand
expression before taking absolute values is
\[
 (m_3-\mu)+\mu J+2s\mu^2-4s^2\mu+2s^3.
\]
Since $|m_3-\mu|\le J$ and $|\mu|,|s|\le1$, its absolute value is at most
$2J+2\mu^2+6s^2$.  If $Z\ge1/2$, division by $Z$ gives the desired
bound.  If $Z<1/2$, use
$Z\ge1-J-\mu^2$, so $J+\mu^2>1/2$, and the trivial bound
$|\mu(F_uw)-\mu(w)+2s|\le4\le8(J+\mu^2)$.
All states used here are regular almost surely.

Let $\Pi_j$ be the normalized phase-$j$ conditional probability and set
\[
 M_j=\mathbb E_{\Pi_j}\mu(w),\qquad
 S_j=\mathbb E_{\Pi_j}s.
\]
All phase-conditioned versions of the preceding nonnegative moment bounds
remain $O_{\bm\sigma}(C^2)$, since the phase probabilities equal $1/L$.
Because $T_\#\Pi_j=\Pi_{j+1}$, \eqref{eq:r01-mean-remainder} gives
\begin{equation}\label{eq:r01-phase-update}
 M_{j+1}=M_j-2S_j+O_{\bm\sigma}(C^2).
\end{equation}
The previous spectral coordinate at phase $j$ has the first-coordinate
law at phase $j-1$.  Also,
\[
 s_1(w)-\mu(w)-r
 =-\frac{r\bigl(J(w)+\mu(w)^2+r\mu(w)\bigr)}{1+r\mu(w)}.
\]
Its conditional $L^1$ norm is $O_{\bm\sigma}(C^3)$ by the established
moment bounds.  Thus
\begin{equation}\label{eq:r01-phase-selector}
 S_j=M_{j-1}+r\sigma_j+O_{\bm\sigma}(C^3).
\end{equation}
Combining these equations and dividing by $r\asymp C$ yields, for
$y_j=M_j/r$,
\[
 y_{j+1}-y_j+2y_{j-1}
 =-2\sigma_j+O_{\bm\sigma}(C).
\]
The cyclic operator on the left is invertible: a nonzero homogeneous
solution has characteristic roots $(1\pm i\sqrt7)/2$, both of modulus
$\sqrt2$, and hence cannot be $L$-periodic.  Its inverse is fixed with the
word.  Since $(-a_j)$ is the exact forced periodic solution,
\[
 y_j=-a_j+O_{\bm\sigma}(C),\qquad
 S_j=C(\sigma_j-a_{j-1})+O_{\bm\sigma}(C^2)
    =C(r_j-1)+O_{\bm\sigma}(C^2).
\]
Finally Jensen's inequality within each phase and the endpoint defect give
\begin{align*}
 \Delta
 &\ge\frac{(1-C)^2}{2L}\sum_j S_j^2\\
 &\ge\frac{V_{\bm\sigma}}2C^2-O_{\bm\sigma}(C^3).
\end{align*}
Together with $G\ge\Delta$ and
\eqref{eq:r01-basic-smallness}, this proves
\[
 G=\frac{V_{\bm\sigma}}2C^2+O_{\bm\sigma}(C^3).
\]
Since $\Gamma_{\bm\sigma}=Ce^{-G}$,
\[
 C-\Gamma_{\bm\sigma}
 =C(1-e^{-G})
 =\frac{V_{\bm\sigma}}2C^3+O_{\bm\sigma}(C^4).
\]
The endpoint trajectory is realizable in dimension two with a positive
warm-up.  Therefore
$\rho_{\bm\sigma}\le q^{\rm fd}_{\bm\sigma}\le\Gamma_{\bm\sigma}$,
and \eqref{eq:r01-endpoint-loggap} gives the same expansion for
$q^{\rm fd}_{\bm\sigma}$.  The identities
$q^{\rm H}_{\bm\sigma}=\Theta^{\rm fd}_{\bm\sigma}=\Gamma_{\bm\sigma}$
give the remaining assertions.  For $01$, the exponents are
$(r_0,r_1)=(3/4,5/4)$, so $V_{01}=1/16$.
\end{proof}


\subsection{How the word controls the coefficient}
% Self-contained R01 mathematical insertion; no unrestricted-factor theorem
% is assumed except in the explicitly conditional paragraph below.
\begin{proposition}[Frequency response of the endpoint coefficient]
\label{prop:r01-frequency}
Let $\bm\sigma\in\{0,1\}^L$, and let $a_n$ be the unique periodic
solution of $a_{n+1}-a_n+2a_{n-1}=2\sigma_n$. Put
$r_n=1-a_{n-1}+\sigma_n$ and
\[
 \widehat\sigma_j=\frac1L\sum_{n=0}^{L-1}\sigma_n e^{-2\pi ijn/L},
 \qquad \omega_j=\frac{2\pi j}{L}.
\]
Then
\begin{equation}\label{eq:r01-parseval}
 \operatorname{Var}(r)=
 \sum_{j=1}^{L-1}
 \frac{1-\cos\omega_j}{1-3\cos\omega_j+4\cos^2\omega_j}
 |\widehat\sigma_j|^2.
\end{equation}
The endpoint cubic coefficient is
$K_{\bm\sigma}:=\operatorname{Var}(r)/2$, that is,
$C-\rho_{\bm\sigma}=K_{\bm\sigma}C^3+O(C^5)$ for each fixed word.
If $t$ is the number of cyclic $1\to0$ boundaries, then
\begin{equation}\label{eq:r01-density-bounds}
 \frac{t}{16L}\le K_{\bm\sigma}\le\frac{8t}{7L}.
\end{equation}
For a nonconstant binary word, equality in the lower bound holds exactly
for alternation, allowing repetitions.
\end{proposition}
\begin{proof}
Write $z=e^{i\omega_j}$ and $q_n=r_n-1$. Summation of the recurrence
gives $\widehat a_0=\widehat\sigma_0$, so $\widehat q_0=0$.
The nonzero modes satisfy
\[
 (z-1+2z^{-1})\widehat a_j=2\widehat\sigma_j,
 \qquad
 \widehat q_j=\frac{z-1}{z-1+2z^{-1}}\widehat\sigma_j.
\]
The denominator never vanishes on the unit circle. Indeed, with
$x=\cos\omega_j$,
\[
 |z-1+2z^{-1}|^2=8x^2-6x+2
 =8(x-3/8)^2+7/8.
\]
Also $|z-1|^2=2(1-x)$. Parseval's identity for the stated $1/L$
normalization proves \eqref{eq:r01-parseval}. The cubic formula follows
by expanding the exact endpoint product, as in the endpoint-response result.

Let $d_n=\sigma_{n+1}-\sigma_n$. Then
$\widehat q_j=\widehat d_j/(z-1+2z^{-1})$ and
$L^{-1}\sum_n|d_n|^2=2t/L$. Since
\[
 7/8\le |z-1+2z^{-1}|^2\le16,
\]
another application of Parseval gives \eqref{eq:r01-density-bounds}.
The maximum 16 is attained on the unit circle only at $z=-1$.
Thus lower-bound equality for a nonconstant word requires all nonzero
Fourier mass of $\bm\sigma$ to be at frequency $\pi$.
A binary sequence with this property alternates between zero and one.
Conversely, its coefficient is $1/32=t/(16L)$.
\end{proof}

Transition frequency controls the size of a guaranteed gap, but does not
determine its exact coefficient. The following examples hold the word length
and the number of each selector fixed:
\begin{center}
\begin{tabular}{lrrr}
\hline
Word & $L$ & $t$ & $K_{\bm\sigma}$\\
\hline
$00010111$ & 8 & 2 & $23/384$\\
$00100111$ & 8 & 2 & $1/12$\\
$000111$   & 6 & 1 & $11/96$\\
$001011$   & 6 & 2 & $15/224$\\
\hline
\end{tabular}
\end{center}
Every displayed word has equal numbers of zeros and ones and is primitive.
For a direct exact check, the centered responses $q=r-1$ are respectively
\[
\begin{aligned}
 &\tfrac1{24}(3,1,-5,17,3,-7,-13,1),\\
 &\tfrac1{3}(0,-2,1,2,0,-1,-1,1),\\
 &\tfrac1{4}(1,3,1,-1,-3,-1),\\
 &\tfrac1{28}(5,-1,17,-9,-15,3).
\end{aligned}
\]
The first two have identical transition density but distinct coefficients.
The last pair increases $t$ from one to two while decreasing the coefficient
by $1/21$; therefore the word with more switches has the larger, slower
endpoint factor for all sufficiently small positive $C$.

\paragraph{Unrestricted interpretation.}
Theorem~\ref{thm:r01-exact-cubic} establishes
\[
 C-\Gamma_{\bm\sigma}=K_{\bm\sigma}C^3+O(C^4)
\]
for every fixed word. The examples therefore give
\[
\begin{aligned}
 \Gamma_{00010111}-\Gamma_{00100111}
    &=\tfrac{3}{128}C^3+O(C^4),\\
 \Gamma_{001011}-\Gamma_{000111}
    &=\tfrac1{21}C^3+O(C^4).
\end{aligned}
\]
Equation~\eqref{eq:r01-density-bounds} bounds the exact
unrestricted leading coefficient as well. These are local comparisons
as $C\downarrow0$ for fixed words, not global orderings for all condition
numbers.

\begin{proposition}[Long blocks]\label{prop:r01-blocks}
For the word $0^m1^m$, write $V_m=\operatorname{Var}(r)$ and set
\[
 \zeta=\frac{1+i\sqrt7}{4},\qquad
 \beta=\frac38-\frac{i}{8\sqrt7}.
\]
Then
\begin{equation}\label{eq:r01-block-exact}
 mV_m=\operatorname{Re}\left[
   \beta\frac{1-\zeta^m}{1+\zeta^m}\right],
\end{equation}
and
\begin{equation}\label{eq:r01-block-error}
 \left|mV_m-\frac38\right|
 \le\frac{2}{\sqrt7}\frac{2^{-m/2}}{1-2^{-m/2}}.
\end{equation}
In particular,
\[
 K_{0^m1^m}=\frac{3}{16m}
       +O\!\left(\frac{2^{-m/2}}m\right).
\]
\end{proposition}
\begin{proof}
The transfer function from $d_n=\sigma_{n+1}-\sigma_n$ to $q_n$ is
\[
 \frac1{z-1+2z^{-1}}=\frac{z}{z^2-z+2}
       =\sum_{j\ge1}h_jz^j,
\]
where $h_0=0$, $h_1=1/2$, and
$h_{j+1}=(h_j-h_{j-1})/2$. Both characteristic roots have modulus
$1/\sqrt2$, so the series and all ensuing sums converge absolutely on
the unit circle. Put $E=\sum_{j\ge1}h_j^2$ and
$R_k=\sum_{j\ge1}h_jh_{j+k}$. The recurrence gives
\[
 R_1=\tfrac12E-\tfrac12R_1,\qquad
 E-\tfrac14=\tfrac12E-\tfrac12R_1.
\]
Hence $E=3/8$ and $R_1=1/8$. For $k\ge2$,
$R_k=(R_{k-1}-R_{k-2})/2$, so
\[
 R_k=\operatorname{Re}(\beta\zeta^k),\qquad
 |R_k|\le\frac1{\sqrt7}2^{-k/2}.
\]
For $0^m1^m$, the difference sequence consists of two opposite unit
impulses separated by $m$ sites. Periodizing the impulse response with
period $2m$ and expanding its squared norm gives
\[
 mV_m=E+2\sum_{k\ge1}(-1)^kR_{km}.
\]
The geometric series evaluates this as \eqref{eq:r01-block-exact},
and the bound on $R_k$ proves \eqref{eq:r01-block-error}.
\end{proof}
The long-block assertion describes the coefficients obtained after taking
$C\downarrow0$ at each fixed $m$, followed by $m\to\infty$.
Section S15 strengthens the remainder to be uniform over all words and proves a joint limit under $m\to\infty$ and $mC\to0$.

The companion directory \texttt{word\_analysis/} contains exact rational response calculations for all 745 primitive rotation classes through length 12, the four displayed examples, and long-block checks for $1\le m\le64$. A separately evaluated discrete Fourier formula agrees within $2.23\,10^{-16}$; five deterministic CSV files reproduce byte for byte in a fresh process. These checks supplement the analytic proofs and do not fit unrestricted factors from finite trajectories.
